% Insert after the R1 display and its middle-line explanation.
% W_m already denotes a weight form in R1; use the explicit source name
% W_m^{src} here, with its identification to the W_m of TB.37--39.
The quotient in this statement is the full-source quotient $q$.
At a proper original source $W\subseteq V$, the same correction
$R=s_h+\Theta T$ has the exact class
\[
 [Ru]_W-[s_hu]_W=\eta_W([Tu]_W),\qquad
 \eta_W:V/W\longrightarrow\mathscr B_\Omega/\Theta W,
 \quad[v]\longmapsto[\Theta v]_W.
\tag{R1a}
\]
Injectivity of $\Theta$ proves injectivity of $\eta_W$, and
subtracting the original source functions proves the equality.
Thus every independent coefficient face retains the kernel
$(V/W)^{\mathfrak a}$ before passing to $q$, including its
specified outer leg label and its empty-coefficient endpoint.
The complete quotient and support maps are (D9)--(D13).

For the finite sources
$W_m^{\rm src}=\operatorname{span}\{\phi_*,D\phi_*,\ldots,D^m\phi_*\}$,
which are the spaces denoted $W_m$ in (TB.37)--(TB.39), the
generator has type
\[
\overline D:
 \mathscr B_\Omega/\Theta W_m^{\rm src}\longrightarrow
 \mathscr B_\Omega/\Theta W_{m+1}^{\rm src}.
\tag{R1b}
\]
It is well-defined because $DW_m^{\rm src}\subseteq W_{m+1}^{\rm src}$
and $D\Theta=\Theta D$.  The full primitive (TB.35) belongs to
$W_{m+1}^{\rm src}$, and $R_{m+1}-R_m\in\Theta W_{m+1}^{\rm src}$.
Taking their original quotient classes therefore gives
\[
 \overline D[R_mu]_{W_m^{\rm src}}
       =[R_{m+1}Au]_{W_{m+1}^{\rm src}},\qquad
 [K_{R_m}u]_{W_m^{\rm src}}
       =-[D^{m+1}\phi_*]_{W_m^{\rm src}}\,d_mu.
\tag{R1c}
\]
The second equality follows from the monic $Q_{m+1}$ and
degree-$m$ $Q_m$ in the full source formula (TB.35).
These exact source-transition identities accompany the unchanged
Grams in (R1); the actual lower source term and full two-row
observation remain in (TB.35)--(TB.36).
